\section{Ioffe 时间分布与伪 PDF}

定义部分子分布的基本对象是一个双定域算符的矩阵元，（略去其自旋结构的不重要细节后）一般可写作 $\langle  p | \phi (0)  \phi(z)  | p \rangle $。由于洛伦兹变换下的不变性，它是两个标量的函数：{\it Ioffe 时间} $(pz)$ \cite{Ioffe:1969kf}（将记作 $-\nu$）与间隔 $z^2$
  \begin{align}
  \langle p |   \phi(0) \phi (z)|p \rangle
=  & {\cal M} (-(pz), -z^2)  =  {\cal M} (\nu, -z^2)
\,
 \
\label{lorentz}
\end{align}
（同样，第二个宗量的符号选取使得\mbox{类空间隔 $z$}对应正值）。可以证明 \cite{Radyushkin:2016hsy,Radyushkin:1983wh}，对所有相关的费曼图，它对 $(pz)$ 的傅里叶变换 ${\cal P} (x, -z^2)$ 以 $-1 \leq x \leq 1$ 为支撑，即
   \begin{align}
 {\cal M} (-(pz), -z^2)
&   =
 \int_{-1}^1 dx
 \, e^{-i x (pz) } \,  {\cal P} (x, -z^2)  \   .
  \label{MPD}
\end{align}
在 $x$ 的这一协变定义中，无需假设 $z$ 位于光锥 $z^2=0$ 上，也无需假设 $p$ 是类光的 $p^2=0$。

在光锥 $z^2=0$ 上，我们形式地有 \mbox{${\cal P} (x, 0) =f(x) $}。因此，函数 ${\cal P} (x, -z^2) $ 可视为把 PDF 概念推广到非类光间隔 $z^2$，依照文献 \cite{Radyushkin:2017cyf}，我们称之为{\it 伪 PDF}。着眼于格点应用，我们取沿强子动量 $p=\{E,0,0,P\}$ 方向取向的间隔 $z=\{0,0,0,z_3\}$。

在可重正化理论（包括 QCD）中，函数 ${\cal M} (\nu, -z^2) $ 具有 $\sim \ln (-z^2)$ 型对数奇异性。在深非弹散射（DIS）中，它们表现为相对于光子虚度 $Q^2$ 的对数标度破坏。一个广为流传的说法是，依赖 $Q^2$ 的 DIS 结构函数 $W(x_B,Q^2)$ 探测距离 $\sim 1/Q$ 处的强子结构。对伪 PDF ${\cal P} (x,z_3^2)$ 而言，可以说它们字面上描述距离 $z_3$ 处的强子结构。

正如 DIS 形状因子 $W(x_B,Q^2)$ 用普适的部分子密度 $f(x, Q^2)$ 表出一样，由格点计算得到的伪 PDF 也可通过通常的部分子分布表达。后者由光锥 $z^2=0$ 上的算符定义，即处于一个对数奇异的极限中。在基于算符乘积展开（OPE）的方法中，标准做法是借助某种方案消除这些奇异性。

其中最流行的是基于维数正规化的 $\overline{\rm MS}$ 方案。于是所得 PDF 依赖于重正化标度 $\mu$，因此应把 PDF 写作 $f(x, \mu^2)$。把 $x$ 换成 Ioffe 时间 $\nu$ 便得到函数
 \begin{align}
 {\cal I} (\nu, \mu^2)
&   =
 \int_{-1}^1 dx
 \, e^{i x \nu  } \,  f (x, \mu^2)  \
  \label{ITD}
\end{align}
它由文献 \cite{Braun:1994jq} 引入并在该文中被称为{\it Ioffe 时间分布}。在此语境下，作为伪 PDF 傅里叶变换的函数 ${\cal M}(\nu , -z^2) $ 应称为 Ioffe 时间{\it 伪分布}或{\it 伪 ITD}。

为建立伪 PDF ${\cal P} (x,z_3^2)$ 与 $\overline{\rm MS}$ 部分子密度 $ f(x, \mu^2)$ 之间的关系，可以对定义 ${\cal P} (x,z_3^2)$ 的矩阵元（即伪 ITD）使用非局域光锥 OPE \cite{Anikin:1978tj,Balitsky:1987bk}（另见 \cite{Radyushkin:2017lvu}）。结果
\begin{align}
{\cal M } (\nu, -z^2 ) =&  \sum_i  \int_{-1}^1 dw \, C_i (w, z^2 \mu ^2, \alpha_s) \,  {\cal I}_i (w \nu,\mu^2)
\nn & + {\cal O} (z^2)
  \  ,
\label{OPE}
\end{align}
具有与 DIS 结构函数 $W (x, Q^2)$ 的通常 OPE 类似的结构。此式中，twist-2 系数函数 $C_i$ 按强耦合常数 $\alpha_s$ 展开，而 ${\cal O} (z^2) $ 代表高扭度项。

然而，把 OPE 应用于 QCD 中的伪 ITD 与伪 PDF 时会遇到与规范连接有关的复杂问题。即当 $z$ 离开光锥时，连接会产生线性 $\sim z_3/a$ 与对数 $\sim \ln (1+z_3^2/a^2)$ 的紫外（UV）发散，其中 $a$ 是带长度量纲的 UV 调节体（可以是有限格距）。虽然这些发散在 $z_3=0$ 时消失，但当 $z_3$ 有限时需要额外的 UV 正规化。

幸运的是，这些发散是乘性的 \mbox{\cite{Polyakov:1980ca,Dotsenko:1979wb,Brandt:1981kf,Aoyama:1981ev,Craigie:1980qs}}（另见近期文献 \cite{Ishikawa:2017faj,Ji:2017oey,Green:2017xeu}），并且在比值——约化 Ioffe 时间分布
 \begin{align}
{\mathfrak M} (\nu, z_3^2) \equiv \frac{ {\cal M} (\nu, z_3^2)}{{\cal M} (0, z_3^2)} \   ,
 \label{redm}
\end{align}
（由我们的论文 \cite{Radyushkin:2017cyf} 引入，其动机部分来自这种相消）——中相互抵消。仅存在于比值分子中的剩余 $\ln z_3^2$ 奇异性由非局域光锥 OPE 描述。

如文献 \cite{Orginos:2017kos} 所述，对小类空间隔 $z^2=-z_3^2$ 且在领头对数层次上，约化伪 PDF 通过第二宗量的简单重标度与 $\overline{\rm MS}$ 分布相联系：
\begin{align}
\mu^2  =4e^{-2\gamma_E}/ z_3^2
  \  ,
\label{Ptof}
\end{align}
其中 $\gamma_E$ 是欧拉常数（更详细的讨论将在下文给出）。由于 \mbox{$2e^{-\gamma_E}=1.12$}，这个重标度因子非常接近 1。然而，系数函数中的 $ {\cal O} (\alpha_s) $ 项可能改变这一因子。
